% TOP_PROBLEM: {"id": 2307033, "title": "Research Problems in Function Theory — Problem 7.33", "category_id": 8, "category": {"id": 8, "name": "analysis", "display_name": "Analysis", "description": "Limits, continuity, calculus, and function theory.", "slug": "analysis", "order_index": 8, "created_at": "2026-07-31T15:26:25.671Z"}, "status": "open"}
% FIRST_POSTED: null
% ENTRY_KIND: statement_only
% Imported from: batch05.tex; UnsolvedMath ID 2307033
\documentclass[11pt]{article}
\usepackage{amsmath,amssymb,mathtools}
\usepackage{fontspec,unicode-math}
\setmainfont{Latin Modern Roman}
\setmathfont{Latin Modern Math}
\usepackage{xurl,hyperref}
\newcommand{\sref}[2]{\hyperlink{source:#1:#2}{[S#2]}}
\newcommand{\cataloguescope}[1]{\paragraph{Mathematical statement.}#1}
\begin{document}
% UnsolvedMath ID 2307033; source catalogue code AMR-022-7033
\setcounter{section}{2307032}
\section{Research Problems in Function Theory — Problem 7.33}\label{2307033}
{\small\sffamily UnsolvedMath ID 2307033 · Analysis}
\subsection{Definitions and mathematical statement}

\paragraph{Shared notation.}$\mathbb N=\{1,2,\ldots\}$, and $\mathbb Z,\mathbb Q,\mathbb R,\mathbb C$ denote the integers, rationals, reals and complex numbers. Any other conventions are specified locally.


For distinct real frequencies $\lambda_1,\ldots,\lambda_N$, let
\[P(\theta)=\sum_{j=1}^N e^{i\lambda_j\theta},\qquad
\mu(P)=\inf_{\theta\in\mathbb R}|P(\theta)|.\]
The rational-frequency special case can be converted, by a common denominator and a unit-modulus factor, to a polynomial with coefficients zero or one evaluated on the unit circle.

\paragraph{Statement recorded in the supplied source.}
Determine sharp bounds for $\mu(P)$ as the distinct frequencies vary, in terms of the number of terms $N$. In particular, improve or determine the sharpness of the recorded universal bound $\mu(P)\le\sqrt{N-1}$ for $N\ge2$, including the rational-frequency polynomial case. \sref{2307033}{2}
The update reports an explicit four-term calculation and examples with few terms and large minimum modulus. \sref{2307033}{2}
\subsection{Short English statement}
Determine how large the minimum modulus of a sum of distinct unit oscillations can be for a given number of terms, including the corresponding zero-one polynomial problem.
\subsection{Sources}
\begin{description}
\item[\hypertarget{source:2307033:1}{S1}] ULAM.AI and the UnsolvedMath Contributors, \emph{UnsolvedMath: A Curated Collection of Open Mathematics Problems}, record 2307033 (AMR-022-7033). CC BY4.0. \url{https://huggingface.co/datasets/ulamai/UnsolvedMath}.
\item[\hypertarget{source:2307033:2}{S2}] W. K. Hayman and E. F. Lingham, \emph{Research Problems in Function Theory} (2018), Chapter7, Problem7.33 and its complete update. \url{https://arxiv.org/abs/1809.07200}.
\end{description}
\label{end:2307033}
\end{document}
